% Options for packages loaded elsewhere
\PassOptionsToPackage{unicode}{hyperref}
\PassOptionsToPackage{hyphens}{url}
\documentclass[
]{ltjsarticle}
\usepackage{xcolor}
\usepackage[margin=2cm]{geometry}
\usepackage{amsmath,amssymb}
\setcounter{secnumdepth}{-\maxdimen} % remove section numbering
\usepackage{iftex}
\ifPDFTeX
  \usepackage[T1]{fontenc}
  \usepackage[utf8]{inputenc}
  \usepackage{textcomp} % provide euro and other symbols
\else % if luatex or xetex
  \usepackage{unicode-math} % this also loads fontspec
  \defaultfontfeatures{Scale=MatchLowercase}
  \defaultfontfeatures[\rmfamily]{Ligatures=TeX,Scale=1}
\fi
\usepackage{lmodern}
\ifPDFTeX\else
  % xetex/luatex font selection
  \setmainfont[]{Hiragino Sans}
  \setmonofont[]{Menlo}
\fi
% Use upquote if available, for straight quotes in verbatim environments
\IfFileExists{upquote.sty}{\usepackage{upquote}}{}
\IfFileExists{microtype.sty}{% use microtype if available
  \usepackage[]{microtype}
  \UseMicrotypeSet[protrusion]{basicmath} % disable protrusion for tt fonts
}{}
\makeatletter
\@ifundefined{KOMAClassName}{% if non-KOMA class
  \IfFileExists{parskip.sty}{%
    \usepackage{parskip}
  }{% else
    \setlength{\parindent}{0pt}
    \setlength{\parskip}{6pt plus 2pt minus 1pt}}
}{% if KOMA class
  \KOMAoptions{parskip=half}}
\makeatother
\usepackage{color}
\usepackage{fancyvrb}
\newcommand{\VerbBar}{|}
\newcommand{\VERB}{\Verb[commandchars=\\\{\}]}
\DefineVerbatimEnvironment{Highlighting}{Verbatim}{commandchars=\\\{\}}
% Add ',fontsize=\small' for more characters per line
\newenvironment{Shaded}{}{}
\newcommand{\AlertTok}[1]{\textcolor[rgb]{1.00,0.00,0.00}{\textbf{#1}}}
\newcommand{\AnnotationTok}[1]{\textcolor[rgb]{0.38,0.63,0.69}{\textbf{\textit{#1}}}}
\newcommand{\AttributeTok}[1]{\textcolor[rgb]{0.49,0.56,0.16}{#1}}
\newcommand{\BaseNTok}[1]{\textcolor[rgb]{0.25,0.63,0.44}{#1}}
\newcommand{\BuiltInTok}[1]{\textcolor[rgb]{0.00,0.50,0.00}{#1}}
\newcommand{\CharTok}[1]{\textcolor[rgb]{0.25,0.44,0.63}{#1}}
\newcommand{\CommentTok}[1]{\textcolor[rgb]{0.38,0.63,0.69}{\textit{#1}}}
\newcommand{\CommentVarTok}[1]{\textcolor[rgb]{0.38,0.63,0.69}{\textbf{\textit{#1}}}}
\newcommand{\ConstantTok}[1]{\textcolor[rgb]{0.53,0.00,0.00}{#1}}
\newcommand{\ControlFlowTok}[1]{\textcolor[rgb]{0.00,0.44,0.13}{\textbf{#1}}}
\newcommand{\DataTypeTok}[1]{\textcolor[rgb]{0.56,0.13,0.00}{#1}}
\newcommand{\DecValTok}[1]{\textcolor[rgb]{0.25,0.63,0.44}{#1}}
\newcommand{\DocumentationTok}[1]{\textcolor[rgb]{0.73,0.13,0.13}{\textit{#1}}}
\newcommand{\ErrorTok}[1]{\textcolor[rgb]{1.00,0.00,0.00}{\textbf{#1}}}
\newcommand{\ExtensionTok}[1]{#1}
\newcommand{\FloatTok}[1]{\textcolor[rgb]{0.25,0.63,0.44}{#1}}
\newcommand{\FunctionTok}[1]{\textcolor[rgb]{0.02,0.16,0.49}{#1}}
\newcommand{\ImportTok}[1]{\textcolor[rgb]{0.00,0.50,0.00}{\textbf{#1}}}
\newcommand{\InformationTok}[1]{\textcolor[rgb]{0.38,0.63,0.69}{\textbf{\textit{#1}}}}
\newcommand{\KeywordTok}[1]{\textcolor[rgb]{0.00,0.44,0.13}{\textbf{#1}}}
\newcommand{\NormalTok}[1]{#1}
\newcommand{\OperatorTok}[1]{\textcolor[rgb]{0.40,0.40,0.40}{#1}}
\newcommand{\OtherTok}[1]{\textcolor[rgb]{0.00,0.44,0.13}{#1}}
\newcommand{\PreprocessorTok}[1]{\textcolor[rgb]{0.74,0.48,0.00}{#1}}
\newcommand{\RegionMarkerTok}[1]{#1}
\newcommand{\SpecialCharTok}[1]{\textcolor[rgb]{0.25,0.44,0.63}{#1}}
\newcommand{\SpecialStringTok}[1]{\textcolor[rgb]{0.73,0.40,0.53}{#1}}
\newcommand{\StringTok}[1]{\textcolor[rgb]{0.25,0.44,0.63}{#1}}
\newcommand{\VariableTok}[1]{\textcolor[rgb]{0.10,0.09,0.49}{#1}}
\newcommand{\VerbatimStringTok}[1]{\textcolor[rgb]{0.25,0.44,0.63}{#1}}
\newcommand{\WarningTok}[1]{\textcolor[rgb]{0.38,0.63,0.69}{\textbf{\textit{#1}}}}
\usepackage{longtable,booktabs,array}
\newcounter{none} % for unnumbered tables
\usepackage{calc} % for calculating minipage widths
% Correct order of tables after \paragraph or \subparagraph
\usepackage{etoolbox}
\makeatletter
\patchcmd\longtable{\par}{\if@noskipsec\mbox{}\fi\par}{}{}
\makeatother
% Allow footnotes in longtable head/foot
\IfFileExists{footnotehyper.sty}{\usepackage{footnotehyper}}{\usepackage{footnote}}
\makesavenoteenv{longtable}
\setlength{\emergencystretch}{3em} % prevent overfull lines
\providecommand{\tightlist}{%
  \setlength{\itemsep}{0pt}\setlength{\parskip}{0pt}}
\usepackage{bookmark}
\IfFileExists{xurl.sty}{\usepackage{xurl}}{} % add URL line breaks if available
\urlstyle{same}
\hypersetup{
  hidelinks,
  pdfcreator={LaTeX via pandoc}}

\author{}
\date{}

\begin{document}

{
\setcounter{tocdepth}{2}
\tableofcontents
}
\section{PhiLang --- Minimal
Interpreter}\label{philang-minimal-interpreter}

進化計算 (本ディレクトリの REPORT.md §6) が示唆した \textbf{Phase
Orchestration Language} の最小実装。

\subsection{構成}\label{ux69cbux6210}

{\def\LTcaptype{none} % do not increment counter
\begin{longtable}[]{@{}llr@{}}
\toprule\noalign{}
ファイル & 役割 & 行数 \\
\midrule\noalign{}
\endhead
\bottomrule\noalign{}
\endlastfoot
\texttt{philang\_parser.py} & \texttt{.phi} → \texttt{Phase} AST &
188 \\
\texttt{philang\_interp.py} & AST 実行 + guarantees 検証 & 207 \\
\texttt{philang\_stdlib.py} & 数値カーネル (現状 Chudnovsky のみ) &
60 \\
\texttt{philang\_main.py} & CLI エントリ & 26 \\
\texttt{compute\_pi\_10k.phi} & サンプル: π 10,000 桁 & 26 \\
\end{longtable}
}

\subsection{実行}\label{ux5b9fux884c}

\begin{Shaded}
\begin{Highlighting}[]
\ExtensionTok{/usr/bin/python3}\NormalTok{ philang\_main.py compute\_pi\_10k.phi}
\end{Highlighting}
\end{Shaded}

期待出力:

\begin{verbatim}
[phi] running phase 'ComputePi10k'
[phi]   compute: chudnovsky_pi(digits=10000, ckpt_every=1000, ckpt_path='.../pi_10k.partial.txt')
[phi]   wrote 10002 chars → .../pi_10k.txt (1.832s)
[phi] verifying guarantees...
[phi]   digit_count: ✓ (10000 ≥ 10000)
[phi]   starts_with: ✓ (37 char prefix matches)
[phi]   no_silent_death: ✓ (reached verification phase)
[phi]   bounded_wall: ✓ (1.83s ≤ 60.0s)
[phi]   partial_resume_on_kill: ✓ (1 checkpoint clause(s))
[phi] phase 'ComputePi10k': all guarantees satisfied in 1.83s
\end{verbatim}

\subsection{言語の意味論}\label{ux8a00ux8a9eux306eux610fux5473ux8ad6}

PhiLang プログラムは \textbf{1 つの phase + 1 つの guarantees
ブロック}で構成される。

\begin{verbatim}
phase NAME
  within   <duration> and <cost>     # ハード予算
  drains   for <dur> idle <dur>      # actor drain buffer (LLM 用、kernel では未使用)
  caffeine if <cond>                 # Mac sleep 対策
  uses     <provider/model>          # 任意
{
  <key> = <value>                    # 通常パラメタ
  every call has policy { ... }      # LLM コール用 (Phase-1 等)
  after every <N> <unit> { ... }     # チェックポイント / 進捗フック
}
guarantees {
  <predicate>                        # bool 形式
  <predicate> = <value>              # 値付き形式
}
\end{verbatim}

\subsubsection{実装済の compute
kernel}\label{ux5b9fux88c5ux6e08ux306e-compute-kernel}

\texttt{philang\_stdlib.COMPUTE\_REGISTRY} で参照可能なもの:

{\def\LTcaptype{none} % do not increment counter
\begin{longtable}[]{@{}
  >{\raggedright\arraybackslash}p{(\linewidth - 4\tabcolsep) * \real{0.3333}}
  >{\raggedright\arraybackslash}p{(\linewidth - 4\tabcolsep) * \real{0.3333}}
  >{\raggedright\arraybackslash}p{(\linewidth - 4\tabcolsep) * \real{0.3333}}@{}}
\toprule\noalign{}
\begin{minipage}[b]{\linewidth}\raggedright
key
\end{minipage} & \begin{minipage}[b]{\linewidth}\raggedright
関数
\end{minipage} & \begin{minipage}[b]{\linewidth}\raggedright
用途
\end{minipage} \\
\midrule\noalign{}
\endhead
\bottomrule\noalign{}
\endlastfoot
\texttt{chudnovsky\_binary\_splitting} &
\texttt{chudnovsky\_pi(digits,\ ckpt\_every,\ ckpt\_path)} & π
を任意桁計算 \\
\end{longtable}
}

\subsubsection{実装済の
guarantee}\label{ux5b9fux88c5ux6e08ux306e-guarantee}

{\def\LTcaptype{none} % do not increment counter
\begin{longtable}[]{@{}
  >{\raggedright\arraybackslash}p{(\linewidth - 2\tabcolsep) * \real{0.5000}}
  >{\raggedright\arraybackslash}p{(\linewidth - 2\tabcolsep) * \real{0.5000}}@{}}
\toprule\noalign{}
\begin{minipage}[b]{\linewidth}\raggedright
名前
\end{minipage} & \begin{minipage}[b]{\linewidth}\raggedright
検証内容
\end{minipage} \\
\midrule\noalign{}
\endhead
\bottomrule\noalign{}
\endlastfoot
\texttt{no\_silent\_death} & guarantees ブロックに到達できたか (=
プロセス生存) \\
\texttt{bounded\_wall\ =\ \textless{}dur\textgreater{}} &
実時間が指定値以下か \\
\texttt{bounded\_cost\ =\ \textless{}cost\textgreater{}} & (LLM
用、kernel は常に 0) \\
\texttt{digit\_count\ =\ \textless{}N\textgreater{}} & 出力の桁数 ≥ N \\
\texttt{starts\_with\ =\ "\textless{}prefix\textgreater{}"} &
出力が指定の文字列で始まるか \\
\texttt{partial\_resume\_on\_kill} &
\texttt{after\ every\ ...\ checkpoint} 句が phase 内にあるか \\
\end{longtable}
}

\subsection{検証実績
(2026-05-16)}\label{ux691cux8a3cux5b9fux7e3e-2026-05-16}

{\def\LTcaptype{none} % do not increment counter
\begin{longtable}[]{@{}
  >{\raggedright\arraybackslash}p{(\linewidth - 2\tabcolsep) * \real{0.5000}}
  >{\raggedright\arraybackslash}p{(\linewidth - 2\tabcolsep) * \real{0.5000}}@{}}
\toprule\noalign{}
\begin{minipage}[b]{\linewidth}\raggedright
項目
\end{minipage} & \begin{minipage}[b]{\linewidth}\raggedright
値
\end{minipage} \\
\midrule\noalign{}
\endhead
\bottomrule\noalign{}
\endlastfoot
計算桁数 & 10,000 \\
所要時間 & 1.83 秒 \\
先頭 102 桁の既知 π との一致 & ✓ \\
Feynman 点 (\texttt{999999} の最初の出現) & 位置 762 (既知値と一致) \\
1000 桁目 & \texttt{9} (既知 OEIS A000796 と一致) \\
出力ファイル & \texttt{pi\_10k.txt} (10003 bytes),
\texttt{pi\_10k.partial.txt} (中間チェックポイント) \\
\end{longtable}
}

\subsection{仕様外 /
既知の限界}\label{ux4ed5ux69d8ux5916-ux65e2ux77e5ux306eux9650ux754c}

\begin{itemize}
\tightlist
\item
  \textbf{\texttt{uses} 句}: 現状は構文上受け取るだけで実際の provider
  切替は行わない (kernel が local なので)。
\item
  \textbf{\texttt{every\ call\ has\ policy}}: LLM コール用の宣言だが、現
  stdlib に LLM kernel が無いので未テスト。Phase-1 等の AIPL
  連携を将来追加するときに使う。
\item
  \textbf{エラー時の partial\_resume\_on\_kill}:
  \texttt{after\ every\ ...\ checkpoint} 句があれば guarantee は OK
  だが、実際に途中で kill されたときの resume ロジックは未実装。
\item
  \textbf{並列実行}: なし。Chudnovsky kernel は単一スレッドで decimal
  を回す。
\end{itemize}

\end{document}
